\section{Min-max normalizace dat}\label{sec:ai-minmax}

Standardizace popisuje hodnoty pomocí vzdálenosti od průměru. Jinou možností je převést nejmenší hodnotu na \(0\), největší na \(1\) a~ostatní mezi ně lineárně rozmístit.

\begin{definition}\label{def:ai-minmax}
    Nechť \(\mathbf{X}=(x_1,\ldots,x_n)\) a
    \[
        m=\min_{1\leqslant i\leqslant n}x_i,
        \qquad
        M=\max_{1\leqslant i\leqslant n}x_i.
    \]
    Je-li \(M>m\), definujeme \emph{min-max normalizované hodnoty}
    \[
        u_i=\frac{x_i-m}{M-m}
    \]
    a~normalizovaný soubor \(\mathbf{U}=(u_1,\ldots,u_n)\).
\end{definition}

Pro každé \(i\) platí \(0\leqslant u_i\leqslant 1\). Nejmenší původní hodnota se zobrazí na \(0\), největší na \(1\) a~pořadí hodnot se nezmění. Jde totiž o~rostoucí lineární transformaci.

\begin{figure}[ht]
    \centering
    \begin{tikzpicture}[x=1cm,y=1cm,every node/.style={font=\normalsize}]
    \draw[diredge] (0,1.5) -- (10,1.5);
    \node[left] at (0,1.5) {\(\mathbf{X}\)};
    \foreach \x/\lab in {1/\(m\),2.6/\(x_2\),4.5/\(x_3\),7/\(x_4\),9/\(M\)}{
        \fill[blue!75!black] (\x,1.5) circle (2.3pt);
        \node[above] at (\x,1.62) {\lab};
    }
    \draw[diredge] (5,0.95) -- node[right]
        {\(u_i=(x_i-m)/(M-m)\)} (5,0.25);
    \draw[diredge] (0,-0.3) -- (10,-0.3);
    \node[left] at (0,-0.3) {\(\mathbf{U}\)};
    \foreach \x/\lab in {1/\(0\),3/\(u_2\),5.4/\(u_3\),7.5/\(u_4\),9/\(1\)}{
        \fill[green!55!black] (\x,-0.3) circle (2.3pt);
        \node[below] at (\x,-0.42) {\lab};
    }
\end{tikzpicture}

    \caption{Min-max normalizace převádí rozsah dat na interval \(\langle0,1\rangle\).}
    \label{fig:ai-minmax}
\end{figure}

Pro soubor teplot
\[
    \mathbf{T}=(18,20,20,21,26)
\]
je \(m=18\), \(M=26\), a~tedy
\[
    \mathbf{U}
    =
    (0;\ 0{,}25;\ 0{,}25;\ 0{,}375;\ 1).
\]

Standardizace i~min-max normalizace odstraňují fyzikální jednotku a~pomáhají porovnávat různé znaky. Liší se však významem výsledných hodnot. Standardizace zachovává informaci o~vzdálenosti od průměru ve směrodatných odchylkách a~není omezená na pevný interval. Min-max normalizace zaručí interval \(\langle0,1\rangle\), ale silně závisí na krajních hodnotách. Pro metody založené na vzdálenosti, například \(k\)-NN nebo \(k\)-means, jsou použitelné obě; vhodnější volba závisí na rozdělení dat a~na výskytu odlehlých hodnot.

\notebox{Stejně jako u~standardizace určujeme \(m\) a~\(M\) pouze z~trénovacích dat. Nová hodnota může ležet mimo původní rozsah, a~její normalizovaná hodnota proto může být menší než \(0\) nebo větší než \(1\).}
